\documentclass[aspectratio=169,11pt]{beamer}

\usepackage[T1]{fontenc}
\usepackage[brazilian]{babel}
\usepackage{booktabs}
\usepackage{tabularx}
\graphicspath{{../../assets/}}
\makeatletter
\IfFileExists{beamerthemeEMASC.sty}{}{\def\input@path{{../modelo/}}}
\makeatother
\usetheme{EMASC}

\setlength{\parskip}{0.35em}
\setlength{\tabcolsep}{6pt}
\renewcommand{\arraystretch}{1.08}
\lstset{style=EMASC,language=Python}
\hypersetup{colorlinks=true,urlcolor=emascPurple,linkcolor=emascPurple,pdfencoding=auto}
\urlstyle{same}
\renewcommand*{\HyperDestNameFilter}[1]{python.#1}

\title{Introdução a Python}
\institute{Liga Acadêmica Matemática \& Empresa}
\date{EMASC 2026}
\AtBeginSection[]{\EMASCsectionpage{INTRODUÇÃO A PYTHON}}

\begin{document}

\EMASCtitlepage{\inserttitle}

\begin{frame}[t]{Roteiro}
\vspace{3mm}
\tableofcontents[hideallsubsections]
\end{frame}

\begin{frame}[t]{Como utilizar esta apostila}
A apostila apresenta os principais fundamentos de Python em uma sequência gradual. Em cada tópico, leia a explicação, observe a estrutura geral do comando e depois execute o exemplo em um interpretador Python 3. Os exemplos podem ser copiados e testados individualmente.\par
\begin{visibleenv}<2->
\begin{EMASCnote}
Observação: Python utiliza indentação (espaços no início das linhas) para delimitar blocos de código. Nos exemplos, mantenha os recuos exatamente como aparecem.
\end{EMASCnote}
\end{visibleenv}
\end{frame}

\section{Variáveis e operadores}

\begin{frame}[fragile,t]{1. Variáveis}
O que são? Variáveis são nomes associados a valores armazenados durante a execução do programa. Em Python, não é necessário declarar previamente o tipo da variável.\par
\begin{columns}[T,onlytextwidth]
\begin{column}{0.46\textwidth}
\EMASCcodelabel{Estrutura do comando}
\begin{lstlisting}
nome_da_variavel = valor
\end{lstlisting}
\end{column}
\begin{column}{0.51\textwidth}
\begin{visibleenv}<2->
\EMASCcodelabel{Exemplo prático}
\begin{lstlisting}
nome = "Carlos"
idade = 20
altura = 1.75

print(nome)
print(idade)
print(altura)
\end{lstlisting}
\end{visibleenv}
\end{column}
\end{columns}
\begin{visibleenv}<3->
\begin{EMASCnote}
Neste exemplo, nome guarda um texto, idade guarda um inteiro e altura guarda um número de ponto flutuante. O sinal \texttt{=} atribui um valor à variável.
\end{EMASCnote}
\end{visibleenv}
\end{frame}

\begin{frame}[t]{2. Tipos de variáveis}
O que são? O tipo indica a categoria do dado armazenado. Os tipos básicos mais comuns são:\par
\begin{tabularx}{\linewidth}{@{}lXl@{}}
\toprule
\textbf{Tipo} & \textbf{Descrição} & \textbf{Exemplo} \\
\midrule
\texttt{int} & Número inteiro & \texttt{idade = 18} \\
\texttt{float} & Número de ponto flutuante & \texttt{altura = 1.75} \\
\texttt{str} & Texto (string) & \texttt{nome = "Ana"} \\
\texttt{bool} & Valor lógico: verdadeiro ou falso & \texttt{aprovado = True} \\
\bottomrule
\end{tabularx}
\par\smallskip
Estrutura: atribua um valor e utilize \texttt{type()} para descobrir seu tipo.\par
\end{frame}

\begin{frame}[fragile,t]{2. Tipos de variáveis / Exemplo}
\begin{columns}[T,onlytextwidth]
\begin{column}{0.46\textwidth}
\EMASCcodelabel{Estrutura do comando}
\begin{lstlisting}
variavel = valor
type(variavel)
\end{lstlisting}
\end{column}
\begin{column}{0.51\textwidth}
\begin{visibleenv}<2->
\EMASCcodelabel{Exemplo prático}
\begin{lstlisting}
idade = 18
altura = 1.70
nome = "João"
aprovado = True

print(type(idade))
print(type(altura))
print(type(nome))
print(type(aprovado))
\end{lstlisting}
\end{visibleenv}
\end{column}
\end{columns}
\begin{visibleenv}<3->
\begin{EMASCnote}
A função \texttt{type()} retorna o tipo do objeto. Os resultados serão, respectivamente, \texttt{int}, \texttt{float}, \texttt{str} e \texttt{bool}.
\end{EMASCnote}
\end{visibleenv}
\end{frame}

\begin{frame}[fragile,t]{3. INPUT e PRINT}
O que é? \texttt{input()} recebe uma informação digitada pelo usuário; \texttt{print()} exibe informações na tela.\par
\begin{columns}[T,onlytextwidth]
\begin{column}{0.40\textwidth}
\EMASCcodelabel{Estrutura do comando}
\begin{lstlisting}
variavel = input("Mensagem: ")
print(variavel)
\end{lstlisting}
\end{column}
\begin{column}{0.56\textwidth}
\begin{visibleenv}<2->
\EMASCcodelabel{Exemplo prático}
\begin{lstlisting}
nome = input("Digite seu nome: ")
idade = int(input("Digite sua idade: "))

print("Seu nome é:", nome)
print("Sua idade é:", idade)
\end{lstlisting}
\end{visibleenv}
\end{column}
\end{columns}
\begin{visibleenv}<3->
\begin{EMASCnote}
Entenda o exemplo: O \texttt{input()} retorna texto (\texttt{str}). Por isso, \texttt{int()} converte a idade para inteiro. Para números de ponto flutuante, utiliza-se \texttt{float()}.
\end{EMASCnote}
\end{visibleenv}
\end{frame}

\begin{frame}[t]{4.1 Operadores aritméticos}
O que são? Operadores realizam cálculos, comparações e combinações lógicas.\par
\begin{tabularx}{\linewidth}{@{}lX@{}}
\toprule
\textbf{Operador} & \textbf{Operação} \\
\midrule
\texttt{+} & Adição \\
\texttt{-} & Subtração \\
\texttt{*} & Multiplicação \\
\texttt{/} & Divisão \\
\texttt{//} & Divisão inteira \\
\texttt{\%} & Resto da divisão \\
\texttt{**} & Potenciação \\
\bottomrule
\end{tabularx}
\end{frame}

\begin{frame}[fragile,t]{4.1 Operadores aritméticos / Exemplo}
\begin{columns}[T,onlytextwidth]
\begin{column}{0.48\textwidth}
\EMASCcodelabel{Estrutura do comando}
\begin{lstlisting}
resultado = valor1 operador valor2
\end{lstlisting}
\end{column}
\begin{column}{0.49\textwidth}
\begin{visibleenv}<2->
\EMASCcodelabel{Exemplo prático}
\begin{lstlisting}
numero1 = 10
numero2 = 5

print(numero1 + numero2) # 15
print(numero1 - numero2) # 5
print(numero1 * numero2) # 50
print(numero1 / numero2) # 2.0
print(numero1 % numero2) # 0
print(numero1 ** numero2) # 100000
\end{lstlisting}
\end{visibleenv}
\end{column}
\end{columns}
\end{frame}

\begin{frame}[fragile,t]{4.2 Operadores relacionais}
Comparam valores e produzem um resultado booleano (\texttt{True} ou \texttt{False}): \texttt{==} igual; \texttt{!=} diferente; \texttt{\textgreater{}} maior; \texttt{\textless{}} menor; \texttt{\textgreater{}=} maior ou igual; \texttt{\textless{}=} menor ou igual.\par
\begin{visibleenv}<2->
\EMASCcodelabel{Exemplo prático}
\begin{lstlisting}
idade = 18

print(idade >= 18)    # True
print(idade == 18)    # True
print(idade != 20)    # True
\end{lstlisting}
\end{visibleenv}
\end{frame}

\begin{frame}[fragile,t]{4.3 Operadores lógicos}
\texttt{and} (E) exige que as condições sejam verdadeiras; \texttt{or} (OU) exige pelo menos uma condição verdadeira; \texttt{not} (NÃO) inverte o valor lógico.\par
\begin{visibleenv}<2->
\EMASCcodelabel{Exemplo prático}
\begin{lstlisting}
idade = 20
tem_carteira = True

print(idade >= 18 and tem_carteira)     # True
\end{lstlisting}
\end{visibleenv}
\end{frame}

\section{Controle e funções}

\begin{frame}[fragile,t]{5. IF, ELIF e ELSE}
O que é? Estruturas condicionais permitem que o programa escolha quais comandos executar de acordo com condições. \texttt{if} testa a primeira condição, \texttt{elif} testa condições adicionais e \texttt{else} executa quando nenhuma condição anterior foi satisfeita.\par
\begin{visibleenv}<2->
\EMASCcodelabel{Estrutura do comando}
\begin{lstlisting}
if condição:
    comando
elif outra_condição:
    comando
else:
    comando
\end{lstlisting}
\end{visibleenv}
\end{frame}

\begin{frame}[fragile,t]{5. IF, ELIF e ELSE / Exemplo}
\EMASCcodelabel{Exemplo prático}
\begin{lstlisting}
nota = float(input("Digite sua nota: "))

if nota >= 7:
    print("Aprovado")
elif nota >= 5:
    print("Recuperação")
else:
    print("Reprovado")
\end{lstlisting}
\begin{visibleenv}<2->
\begin{EMASCnote}
Entenda o exemplo: Se a nota for 7 ou mais, o resultado será aprovado. Se for pelo menos 5 e menor que 7, será recuperação. Abaixo de 5, será reprovado. A ordem das condições importa.
\end{EMASCnote}
\end{visibleenv}
\end{frame}

\begin{frame}[t]{6. WHILE e FOR}
São estruturas de repetição (laços), usadas para executar um bloco de comandos mais de uma vez.\par
\begin{visibleenv}<2->
\begin{EMASCnote}
Diferença prática: use \texttt{while} quando a repetição depende de uma condição; use \texttt{for} quando quiser percorrer uma sequência ou repetir com uma quantidade definida.
\end{EMASCnote}
\end{visibleenv}
\end{frame}

\begin{frame}[fragile,t]{6.1 WHILE}
O \texttt{while} repete os comandos enquanto a condição for verdadeira.\par
\begin{columns}[T,onlytextwidth]
\begin{column}{0.46\textwidth}
\EMASCcodelabel{Estrutura do comando}
\begin{lstlisting}
while condição:
    comando
\end{lstlisting}
\end{column}
\begin{column}{0.51\textwidth}
\begin{visibleenv}<2->
\EMASCcodelabel{Exemplo prático}
\begin{lstlisting}
contador = 1

while contador <= 5:
    print(contador)
    contador += 1
\end{lstlisting}
\end{visibleenv}
\end{column}
\end{columns}
\begin{visibleenv}<3->
\begin{EMASCnote}
O contador começa em 1 e aumenta uma unidade a cada repetição. Quando chega a 6, a condição se torna falsa e o laço termina. Se a variável de controle não for atualizada corretamente, pode ocorrer um laço infinito.
\end{EMASCnote}
\end{visibleenv}
\end{frame}

\begin{frame}[fragile,t]{6.2 FOR}
O \texttt{for} percorre elementos de uma sequência ou repete comandos com base em uma sequência de valores.\par
\EMASCcodelabel{Estrutura do comando}
\begin{lstlisting}
for variavel in sequência:
    comando
\end{lstlisting}
\begin{visibleenv}<2->
\EMASCcodelabel{Exemplo prático}
\begin{lstlisting}
for numero in range(1, 6):
    print(numero)
\end{lstlisting}
\end{visibleenv}
\begin{visibleenv}<3->
\begin{EMASCnote}
\texttt{range(1, 6)} produz os valores 1, 2, 3, 4 e 5. O limite final não é incluído.
\end{EMASCnote}
\end{visibleenv}
\end{frame}

\begin{frame}[fragile,t]{7. Funções}
O que é? Funções são blocos reutilizáveis de código que realizam tarefas específicas. Elas ajudam a organizar o programa e evitar repetição.\par
\begin{columns}[T,onlytextwidth]
\begin{column}{0.46\textwidth}
\EMASCcodelabel{Estrutura do comando}
\begin{lstlisting}
def nome_da_funcao():
    comandos

nome_da_funcao()
\end{lstlisting}
\end{column}
\begin{column}{0.51\textwidth}
\begin{visibleenv}<2->
\EMASCcodelabel{Exemplo prático}
\begin{lstlisting}
def saudacao():
    print("Olá! Seja bem-vindo!")

saudacao()
\end{lstlisting}
\end{visibleenv}
\end{column}
\end{columns}
\begin{visibleenv}<3->
\begin{EMASCnote}
Entenda o exemplo: A palavra \texttt{def} define a função. A linha \texttt{saudacao()} chama a função, fazendo seu bloco ser executado.
\end{EMASCnote}
\end{visibleenv}
\end{frame}

\begin{frame}[fragile,t]{7. Funções com parâmetros e retorno}
\begin{columns}[T,onlytextwidth]
\begin{column}{0.46\textwidth}
\EMASCcodelabel{Estrutura do comando}
\begin{lstlisting}
def nome_da_funcao(parametro):
    return resultado
\end{lstlisting}
\end{column}
\begin{column}{0.51\textwidth}
\begin{visibleenv}<2->
\EMASCcodelabel{Exemplo prático}
\begin{lstlisting}
def somar(a, b):
    resultado = a + b
    return resultado

total = somar(10, 5)
print(total) # 15
\end{lstlisting}
\end{visibleenv}
\end{column}
\end{columns}
\begin{visibleenv}<3->
\begin{EMASCnote}
Os parâmetros a e b recebem os valores enviados na chamada. \texttt{return} devolve o resultado para quem chamou a função.
\end{EMASCnote}
\end{visibleenv}
\end{frame}

\section{Listas, tuplas e dicionários}

\begin{frame}[fragile,t]{8. Listas}
O que é? Listas armazenam vários elementos em uma única estrutura. São mutáveis, isto é, podem ser alteradas depois de criadas, e são representadas por colchetes.\par
\begin{visibleenv}<2->
\EMASCcodelabel{Estrutura do comando}
\begin{lstlisting}
nome_da_lista = [elemento1, elemento2, elemento3]
\end{lstlisting}
\end{visibleenv}
\end{frame}

\begin{frame}[fragile,t]{8. Listas / Exemplo}
\EMASCcodelabel{Exemplo prático}
\begin{lstlisting}
frutas = ["Maçã", "Banana", "Uva"]

print(frutas)
print(frutas[0])   # Maçã
print(frutas[1])   # Banana
\end{lstlisting}
\begin{visibleenv}<2->
\begin{EMASCnote}
Entenda o exemplo: Os índices começam em zero: o índice 0 acessa o primeiro elemento, o índice 1 acessa o segundo e assim por diante.
\end{EMASCnote}
\end{visibleenv}
\end{frame}

\begin{frame}[fragile,t]{8. Listas / Operações comuns}
\EMASCcodelabel{Exemplo prático}
\begin{lstlisting}
frutas.append("Laranja")    # Adiciona ao final
frutas.remove("Banana")     # Remove pelo valor
frutas.pop()                # Remove o último elemento
print(len(frutas))          # Quantidade de elementos
\end{lstlisting}
\end{frame}

\begin{frame}[fragile,t]{9. Tuplas}
O que é? Tuplas armazenam uma sequência de elementos e, ao contrário das listas, são imutáveis: não permitem alterar, adicionar ou remover elementos diretamente após a criação. São normalmente representadas por parênteses.\par
\begin{visibleenv}<2->
\EMASCcodelabel{Estrutura do comando}
\begin{lstlisting}
nome_da_tupla = (elemento1, elemento2, elemento3)
\end{lstlisting}
\end{visibleenv}
\end{frame}

\begin{frame}[fragile,t]{9. Tuplas / Exemplo}
\EMASCcodelabel{Exemplo prático}
\begin{lstlisting}
dias = ("Segunda", "Terça", "Quarta", "Quinta")

print(dias)
print(dias[0])
print(dias[2])
\end{lstlisting}
\begin{visibleenv}<2->
\begin{EMASCnote}
Entenda o exemplo: Assim como nas listas, os índices começam em zero. Uma tupla de um único elemento precisa de vírgula: \texttt{(5,)}.
\end{EMASCnote}
\end{visibleenv}
\end{frame}

\begin{frame}[fragile,t]{10. Dicionários}
O que é? Dicionários armazenam informações em pares de chave e valor. Cada chave identifica o valor associado. São representados por chaves.\par
\begin{visibleenv}<2->
\EMASCcodelabel{Estrutura do comando}
\begin{lstlisting}
nome_do_dicionario = {
    "chave1": valor1,
    "chave2": valor2
}
\end{lstlisting}
\end{visibleenv}
\end{frame}

\begin{frame}[fragile,t]{10. Dicionários / Exemplo}
\EMASCcodelabel{Exemplo prático}
\begin{lstlisting}
aluno = {
    "nome": "Carlos",
    "idade": 20,
    "curso": "Informática"
}

print(aluno["nome"])
print(aluno["idade"])
print(aluno["curso"])
\end{lstlisting}
\begin{visibleenv}<2->
\begin{EMASCnote}
Entenda o exemplo: As chaves são nome, idade e curso. Os valores associados são Carlos, 20 e Informática.
\end{EMASCnote}
\end{visibleenv}
\end{frame}

\begin{frame}[fragile,t]{10. Dicionários / Operações comuns}
\EMASCcodelabel{Exemplo prático}
\begin{lstlisting}
aluno["nota"] = 9.0         # Adiciona uma chave
aluno["idade"] = 21         # Atualiza um valor
aluno.pop("curso")          # Remove a chave e seu valor
print(aluno)
\end{lstlisting}
\end{frame}

\section{Vetores e matrizes}

\begin{frame}[fragile,t]{11. Vetores}
O que é? Em fundamentos de programação, um vetor é uma estrutura unidimensional, organizada em posições sequenciais. Em Python introdutório, uma lista pode representar esse tipo de estrutura.\par
\begin{visibleenv}<2->
\EMASCcodelabel{Estrutura do comando}
\begin{lstlisting}
vetor = [elemento1, elemento2, elemento3]
\end{lstlisting}
\end{visibleenv}
\end{frame}

\begin{frame}[fragile,t]{11. Vetores / Exemplo}
\EMASCcodelabel{Exemplo prático}
\begin{lstlisting}
numeros = [10, 20, 30, 40, 50]

for i in range(len(numeros)):
    print(numeros[i])
\end{lstlisting}
\begin{visibleenv}<2->
\begin{EMASCnote}
Entenda o exemplo: \texttt{len(numeros)} informa a quantidade de elementos. \texttt{range()} gera os índices e \texttt{numeros[i]} acessa o elemento na posição indicada.
\end{EMASCnote}
\end{visibleenv}
\end{frame}

\begin{frame}[fragile,t]{11. Vetores / Acessando posições específicas}
\EMASCcodelabel{Exemplo prático}
\begin{lstlisting}
numeros = [10, 20, 30, 40, 50]

print(numeros[0])   # 10
print(numeros[2])   # 30
print(numeros[4])   # 50
\end{lstlisting}
\end{frame}

\begin{frame}[fragile,t]{12. Matrizes}
O que é? Matrizes organizam dados em linhas e colunas. Uma matriz bidimensional pode ser representada em Python por uma lista cujos elementos também são listas.\par
\begin{visibleenv}<2->
\EMASCcodelabel{Estrutura do comando}
\begin{lstlisting}
matriz = [
    [elemento1, elemento2],
    [elemento3, elemento4]
]
\end{lstlisting}
\end{visibleenv}
\end{frame}

\begin{frame}[fragile,t]{12. Matrizes / Exemplo}
\EMASCcodelabel{Exemplo prático}
\begin{lstlisting}
matriz = [
    [1, 2, 3],
    [4, 5, 6],
    [7, 8, 9]
]

print(matriz[0][0])    # 1
print(matriz[1][2])    # 6
print(matriz[2][1])    # 8
\end{lstlisting}
\begin{visibleenv}<2->
\begin{EMASCnote}
Entenda o exemplo: O primeiro índice indica a linha e o segundo indica a coluna. Como os índices começam em zero, \texttt{matriz[1][2]} acessa a segunda linha e a terceira coluna.
\end{EMASCnote}
\end{visibleenv}
\end{frame}

\begin{frame}[fragile,t]{12. Matrizes / Percorrendo todos os elementos}
\EMASCcodelabel{Exemplo prático}
\begin{lstlisting}
for linha in matriz:
    for elemento in linha:
        print(elemento)
\end{lstlisting}
\begin{visibleenv}<2->
\begin{EMASCnote}
O primeiro \texttt{for} percorre cada linha. O segundo percorre os elementos da linha atual.
\end{EMASCnote}
\end{visibleenv}
\end{frame}

\begin{frame}[t]{13. Resumo geral}
\footnotesize
\renewcommand{\arraystretch}{0.95}
\begin{tabularx}{\linewidth}{@{}lX@{}}
\toprule
\textbf{Assunto} & \textbf{Comando / representação} \\
\midrule
Variáveis & \texttt{nome = valor} \\
Tipos & \texttt{int, float, str, bool} \\
Entrada & \texttt{input()} \\
Saída & \texttt{print()} \\
Operadores & \texttt{+ - * / // \% **; == != \textgreater{} \textless{} \textgreater{}= \textless{}=} \\
Condicionais & \texttt{if, elif, else} \\
Repetição & \texttt{while, for} \\
Funções & \texttt{def, return} \\
Listas & \texttt{[]} \\
Tuplas & \texttt{()} \\
Dicionários & \texttt{\{\} com pares chave: valor} \\
Vetores & lista unidimensional \\
Matrizes & listas aninhadas (linhas e colunas) \\
\bottomrule
\end{tabularx}
\end{frame}

\begin{frame}[t,label=referencias-python]{Referências}
  \small
  \vspace{2mm}
  \textbf{PYTHON SOFTWARE FOUNDATION.} \emph{O tutorial do Python}. Documentação Python 3.\par
  {\footnotesize Disponível em: \url{https://docs.python.org/pt-br/3/tutorial/}.\par}
  \medskip
  \textbf{UNIVERSIDADE FEDERAL DE SANTA CATARINA.} \emph{Funções básicas de entrada e saída de dados}.\par
  {\footnotesize Disponível em: \url{https://mariani.inf.ufsc.br/python/Basico/B04-Entrada_saida.html}.\par}
  \medskip
  \textbf{UNIVERSIDADE FEDERAL DO RIO DE JANEIRO.} Instituto de Computação.
  \emph{Interação com usuário e programa principal}. Python UFRJ.\par
  {\footnotesize Disponível em: \url{https://python.ic.ufrj.br/aulas/aula-10/index.html}.\par}
  \medskip
  {\footnotesize Acesso em: 12 nov. 2026.}
\end{frame}


\end{document}
